\documentclass[10pt]{amsart}
\usepackage{geometry}                % See geometry.pdf to learn the layout options. There are lots.
\geometry{letterpaper}                   % ... or a4paper or a5paper or ... 
%\geometry{landscape}                % Activate for for rotated page geometry
%\usepackage[parfill]{parskip}    % Activate to begin paragraphs with an empty line rather than an indent
\usepackage{graphicx}
\usepackage{amssymb}
\usepackage{epstopdf}
\usepackage{hyperref}
\DeclareGraphicsRule{.tif}{png}{.png}{`convert #1 `dirname #1`/`basename #1 .tif`.png}

\title{Notes on solution to the momentum equation with exponential $\nu$ profile}
\author{GPT 5.5 with  prompting from TWNH }
\date{7Jul26}                                           % Activate to display a given date or no date

\begin{document}
\maketitle

\section{Problem statement I}

Consider the complex horizontal velocity
\begin{equation}
u = u^x+i u^y,
\end{equation}
satisfying
\begin{equation}
E k\,\frac{\partial}{\partial z}
\left(
\nu(z)\frac{\partial u}{\partial z}
\right)
-
i f u
=
\mathfrak f,
\label{eq:main}
\end{equation}
on the vertical interval
\[
-H\le z\le 0.
\]

The boundary conditions are
\begin{align}
E k\,\nu(0)u_z(0) &= \tau^x+i\tau^y, \label{eq:surface_bc}\\
u(-H)&=0. \label{eq:bottom_bc}
\end{align}

Here
\begin{equation}
\mathfrak f
=
p_x+i p_y
\label{eq:pressure_gradient}
\end{equation}
is the complex pressure-gradient forcing and
\begin{equation}
\tau
=
\tau^x+i\tau^y
\label{eq:complex_stress}
\end{equation}
is the complex surface wind stress.

We assume an exponentially varying vertical viscosity
\begin{equation}
\nu(z)
=
a+\beta \exp\!\left[-\frac{z+H}{h}\right],
\label{eq:viscosity}
\end{equation}
with \(a>0\), \(h>0\), and \(\nu(z)>0\) on \([-H,0]\).

The horizontal coordinates \(x,y\) enter parametrically through \(\mathfrak f(x,y,z)\) and \(\tau(x,y)\). In the derivation below this dependence is suppressed.

\section{Reduced form}

Divide \eqref{eq:main} by \(E k\). Then
\begin{equation}
\frac{\partial}{\partial z}
\left(
\nu u_z
\right)
-
\alpha u
=
F,
\label{eq:reduced}
\end{equation}
where
\begin{equation}
\alpha
=
\frac{i f}{E k},
\qquad
F
=
\frac{\mathfrak f}{E k}.
\label{eq:alpha_F}
\end{equation}

Define the differential operator
\begin{equation}
L u
=
\partial_z(\nu u_z)-\alpha u.
\label{eq:operator}
\end{equation}

The problem becomes
\begin{equation}
L u = F,
\label{eq:LuF}
\end{equation}
with boundary conditions
\begin{align}
\nu(0)u_z(0) &= \frac{\tau}{E k}, \label{eq:bc_surface_reduced}\\
u(-H)&=0. \label{eq:bc_bottom_reduced}
\end{align}

\section{Homogeneous equation}

The associated homogeneous problem is
\begin{equation}
\partial_z(\nu u_z)-\alpha u=0.
\label{eq:hom}
\end{equation}

With
\[
\nu(z)=a+\beta \exp\!\left[-\frac{z+H}{h}\right],
\]
introduce
\begin{equation}
X(z)
=
-\frac{\beta}{a}
\exp\!\left[-\frac{z+H}{h}\right].
\label{eq:X_def}
\end{equation}
Then
\begin{equation}
\nu(z)=a(1-X).
\end{equation}

Also define
\begin{equation}
\lambda
=
h\sqrt{\frac{\alpha}{a}}
=
h\sqrt{\frac{i f}{E k\,a}}.
\label{eq:lambda_def}
\end{equation}
Since \(\alpha\) is complex, \(\lambda\) is generally complex. A consistent branch of the square root should be chosen.

In terms of \(X\), the homogeneous equation admits a basis
\begin{align}
y_+(z)
&=
X(z)^\lambda
{}_2F_1
\left(
\lambda,\lambda+1;2\lambda+1;X(z)
\right),
\label{eq:yplus}
\\[0.5em]
y_-(z)
&=
X(z)^{-\lambda}
{}_2F_1
\left(
-\lambda,1-\lambda;1-2\lambda;X(z)
\right).
\label{eq:yminus}
\end{align}

Thus the homogeneous solution is
\begin{equation}
u_h(z)
=
C_+ y_+(z)+C_- y_-(z).
\label{eq:uh}
\end{equation}

\section{Boundary-adapted homogeneous solutions}

We now construct two homogeneous solutions adapted to the boundary conditions.

Let \(\phi(z)\) satisfy
\begin{equation}
L\phi=0,
\qquad
\phi(-H)=0.
\label{eq:phi_def}
\end{equation}
A convenient choice is
\begin{equation}
\phi(z)
=
y_-(-H)y_+(z)-y_+(-H)y_-(z).
\label{eq:phi}
\end{equation}
Indeed,
\[
\phi(-H)
=
y_-(-H)y_+(-H)-y_+(-H)y_-(-H)
=
0.
\]

Next, let \(\psi(z)\) satisfy
\begin{equation}
L\psi=0,
\qquad
\psi_z(0)=0.
\label{eq:psi_def}
\end{equation}
Because the homogeneous surface-stress condition is
\[
\nu(0)\psi_z(0)=0,
\]
and because \(\nu(0)\ne 0\), this is equivalent to \(\psi_z(0)=0\).

A convenient choice is
\begin{equation}
\psi(z)
=
y_-'(0)y_+(z)-y_+'(0)y_-(z).
\label{eq:psi}
\end{equation}
Then
\[
\psi_z(0)
=
y_-'(0)y_+'(0)-y_+'(0)y_-'(0)
=
0.
\]

\section{Weighted Wronskian}

For two homogeneous solutions \(\phi\) and \(\psi\), define the weighted Wronskian
\begin{equation}
W
=
\nu(z)
\left[
\phi(z)\psi_z(z)-\phi_z(z)\psi(z)
\right].
\label{eq:W}
\end{equation}

This quantity is independent of \(z\). Indeed,
\begin{align}
\frac{dW}{dz}
&=
\frac{d}{dz}
\left\{
\nu
\left(
\phi\psi_z-\phi_z\psi
\right)
\right\} \notag\\
&=
\phi(\nu\psi_z)_z-\psi(\nu\phi_z)_z \notag\\
&=
\phi(\alpha\psi)-\psi(\alpha\phi) \notag\\
&=0.
\end{align}

\section{Green's function}

Define \(G(z,\zeta)\) by
\begin{equation}
L_z G(z,\zeta)=\delta(z-\zeta),
\label{eq:green_def}
\end{equation}
with homogeneous boundary conditions
\begin{align}
G(-H,\zeta)&=0, \label{eq:G_bottom}\\
G_z(0,\zeta)&=0. \label{eq:G_surface}
\end{align}

Then the Green's function is
\begin{equation}
G(z,\zeta)
=
\begin{cases}
\dfrac{\phi(z)\psi(\zeta)}{W},
& z<\zeta,\\[1.2em]
\dfrac{\phi(\zeta)\psi(z)}{W},
& z>\zeta.
\end{cases}
\label{eq:green}
\end{equation}

At \(z=\zeta\), \(G\) is continuous:
\begin{equation}
G(\zeta^-,\zeta)=G(\zeta^+,\zeta),
\end{equation}
and its derivative has jump
\begin{equation}
\nu(\zeta)
\left[
G_z(\zeta^+,\zeta)-G_z(\zeta^-,\zeta)
\right]
=
1.
\label{eq:jump}
\end{equation}

\section{Treatment of the nonhomogeneous surface stress}

The Green's function above satisfies homogeneous surface flux condition \(G_z(0,\zeta)=0\). Therefore the pressure-gradient forcing contributes
\begin{equation}
u_{\mathrm{int}}(z)
=
\int_{-H}^{0}
G(z,\zeta)F(\zeta)\,d\zeta
=
\frac{1}{E k}
\int_{-H}^{0}
G(z,\zeta)\mathfrak f(\zeta)\,d\zeta.
\label{eq:uint}
\end{equation}

The nonhomogeneous surface stress is included as a boundary forcing.

Let \(q\) denote the imposed reduced surface flux:
\begin{equation}
q=\nu(0)u_z(0)=\frac{\tau}{E k}.
\label{eq:q}
\end{equation}

Using Green's identity for
\[
L=\partial_z(\nu\partial_z)-\alpha,
\]
one obtains the boundary contribution
\begin{equation}
u_{\mathrm{surf}}(z)
=
- q\,G(z,0).
\label{eq:usurf}
\end{equation}

Hence
\begin{equation}
u_{\mathrm{surf}}(z)
=
-\frac{\tau}{E k}G(z,0).
\label{eq:usurf_tau}
\end{equation}

\section{Full solution}

Combining the internal pressure forcing and the surface-stress forcing gives
\begin{equation}
\boxed{
u(z)
=
\frac{1}{E k}
\int_{-H}^{0}
G(z,\zeta)\mathfrak f(\zeta)\,d\zeta
-
\frac{\tau}{E k}G(z,0).
}
\label{eq:solution_1d}
\end{equation}

Restoring the horizontal dependence,
\begin{equation}
\boxed{
u(x,y,z)
=
\frac{1}{E k}
\int_{-H}^{0}
G(z,\zeta)
\left[
p_x(x,y,\zeta)+i p_y(x,y,\zeta)
\right]
\,d\zeta
-
\frac{\tau^x(x,y)+i\tau^y(x,y)}{E k}
G(z,0).
}
\label{eq:solution_full}
\end{equation}

\section{Summary of definitions}

The viscosity is
\begin{equation}
\nu(z)
=
a+\beta \exp\!\left[-\frac{z+H}{h}\right].
\end{equation}

The transformed variable is
\begin{equation}
X(z)
=
-\frac{\beta}{a}
\exp\!\left[-\frac{z+H}{h}\right],
\end{equation}
so that
\begin{equation}
\nu(z)=a(1-X).
\end{equation}

The complex decay parameter is
\begin{equation}
\lambda
=
h\sqrt{\frac{i f}{E k\,a}}.
\end{equation}

The homogeneous basis functions are
\begin{align}
y_+(z)
&=
X(z)^\lambda
{}_2F_1
\left(
\lambda,\lambda+1;2\lambda+1;X(z)
\right),
\\[0.5em]
y_-(z)
&=
X(z)^{-\lambda}
{}_2F_1
\left(
-\lambda,1-\lambda;1-2\lambda;X(z)
\right).
\end{align}

The boundary-adapted solutions are
\begin{align}
\phi(z)
&=
y_-(-H)y_+(z)-y_+(-H)y_-(z),
\\[0.5em]
\psi(z)
&=
y_-'(0)y_+(z)-y_+'(0)y_-(z).
\end{align}

The weighted Wronskian is
\begin{equation}
W
=
\nu(z)
\left[
\phi(z)\psi_z(z)-\phi_z(z)\psi(z)
\right],
\end{equation}
which is independent of \(z\).

The Green's function is
\begin{equation}
G(z,\zeta)
=
\begin{cases}
\dfrac{\phi(z)\psi(\zeta)}{W},
& z<\zeta,\\[1.2em]
\dfrac{\phi(\zeta)\psi(z)}{W},
& z>\zeta.
\end{cases}
\end{equation}

The final solution is
\begin{equation}
\boxed{
u(x,y,z)
=
\frac{1}{E k}
\int_{-H}^{0}
G(z,\zeta)
\left[
p_x(x,y,\zeta)+i p_y(x,y,\zeta)
\right]d\zeta
-
\frac{\tau^x(x,y)+i\tau^y(x,y)}{E k}
G(z,0).
}
\end{equation}

\section{Constant-viscosity limit}

If \(\beta=0\), then
\[
\nu(z)=a,
\]
and the equation reduces to
\begin{equation}
a u_{zz}-\alpha u=F.
\end{equation}

Define
\begin{equation}
m=\sqrt{\frac{\alpha}{a}}
=
\sqrt{\frac{i f}{E k\,a}}.
\end{equation}

The homogeneous solutions are
\begin{equation}
u_h(z)
=
C_1\cosh\!\left[m(z+H)\right]
+
C_2\sinh\!\left[m(z+H)\right].
\end{equation}

A bottom-Dirichlet solution is
\begin{equation}
\phi(z)
=
\sinh\!\left[m(z+H)\right],
\end{equation}
and a homogeneous surface-stress solution may be chosen as
\begin{equation}
\psi(z)
=
\cosh(mz),
\end{equation}
since
\[
\psi_z(0)=0.
\]

The weighted Wronskian becomes
\begin{equation}
W
=
a
\left[
\phi\psi_z-\phi_z\psi
\right],
\end{equation}
which is constant.

This gives the usual finite-depth Ekman-layer solution with bottom no-slip and imposed surface stress.

\section{Remarks}

\begin{itemize}
    \item The operator
    \[
    L=\partial_z(\nu\partial_z)-\alpha
    \]
    is formally self-adjoint for real \(\alpha\), but here \(\alpha=i f/(E k)\) is complex. The Green's-function construction nevertheless applies for the complex-valued linear ordinary differential equation.

    \item The branch of
    \[
    \lambda=h\sqrt{\frac{i f}{E k\,a}}
    \]
    should be chosen consistently. For \(f>0\), the principal square root gives
    \[
    \sqrt{i}=\frac{1+i}{\sqrt{2}}.
    \]

    \item The condition \(\nu(z)>0\) requires
    \[
    a+\beta \exp\!\left[-\frac{z+H}{h}\right]>0
    \]
    for all \(z\in[-H,0]\).

    \item If \(\beta>0\), then \(X(z)<0\). The hypergeometric expressions should be interpreted using a consistent branch for \(X^\lambda\).

    \item The horizontal coordinates enter only parametrically through the pressure-gradient forcing and the wind stress.
\end{itemize}

\begin{thebibliography}{9}

\bibitem{Ekman1905}
V. W. Ekman.
\newblock On the influence of the Earth's rotation on ocean-currents.
\newblock \emph{Arkiv f\"or Matematik, Astronomi och Fysik}, 2(11):1--52, 1905.

\bibitem{Pedlosky1987}
J. Pedlosky.
\newblock \emph{Geophysical Fluid Dynamics}.
\newblock 2nd edition, Springer, 1987.

\bibitem{Vallis2017}
G. K. Vallis.
\newblock \emph{Atmospheric and Oceanic Fluid Dynamics}.
\newblock 2nd edition, Cambridge University Press, 2017.

\bibitem{Gill1982}
A. E. Gill.
\newblock \emph{Atmosphere--Ocean Dynamics}.
\newblock Academic Press, 1982.

\bibitem{DLMF}
NIST Digital Library of Mathematical Functions.
\newblock \emph{Gauss Hypergeometric Function}.
\newblock Sections 15.2--15.10.
\newblock \url{https://dlmf.nist.gov/15}

\bibitem{Olver2010}
F. W. J. Olver, D. W. Lozier, R. F. Boisvert, and C. W. Clark, editors.
\newblock \emph{NIST Handbook of Mathematical Functions}.
\newblock Cambridge University Press, 2010.

\bibitem{AbramowitzStegun}
M. Abramowitz and I. A. Stegun.
\newblock \emph{Handbook of Mathematical Functions with Formulas, Graphs, and Mathematical Tables}.
\newblock Dover Publications, 1965.

\bibitem{Zettl}
A. Zettl.
\newblock \emph{Sturm--Liouville Theory}.
\newblock American Mathematical Society, 2005.

\bibitem{CoddingtonLevinson}
E. A. Coddington and N. Levinson.
\newblock \emph{Theory of Ordinary Differential Equations}.
\newblock McGraw--Hill, 1955.

\end{thebibliography}

\clearpage

\section{Problem statement II}

We consider the complex horizontal velocity
\begin{equation}
u = u^x+i u^y,
\end{equation}
satisfying the vertical boundary-value problem
\begin{equation}
E k\,\frac{\partial}{\partial z}
\left(
\nu(z)\frac{\partial u}{\partial z}
\right)
-
i f u
=
\mathfrak f,
\label{eq:main}
\end{equation}
on
\[
-H\le z\le 0.
\]

The boundary conditions are
\begin{align}
E k\,\nu(0)u_z(0) &= \tau^x+i\tau^y, \label{eq:surface_bc}\\
u(-H)&=0. \label{eq:bottom_bc}
\end{align}

The complex pressure-gradient forcing is
\begin{equation}
\mathfrak f
=
p_x+i p_y,
\label{eq:pressure_gradient}
\end{equation}
and the complex surface stress is
\begin{equation}
\tau
=
\tau^x+i\tau^y.
\label{eq:stress}
\end{equation}

We now take the vertical viscosity to be
\begin{equation}
\nu(z)
=
a+b\exp\!\left[-\frac{z+H}{h}\right]
+
c\exp\!\left(\frac{z}{h}\right),
\label{eq:viscosity_z}
\end{equation}
where \(a,b,c,h\) are constants with \(h>0\). We assume
\begin{equation}
\nu(z)>0
\qquad
\text{for all } z\in[-H,0].
\end{equation}

The horizontal variables \(x,y\) enter only parametrically through
\[
\mathfrak f(x,y,z)
\quad\text{and}\quad
\tau(x,y).
\]

\section{Reduced form}

Divide \eqref{eq:main} by \(E k\). Define
\begin{equation}
\alpha
=
\frac{i f}{E k},
\qquad
F
=
\frac{\mathfrak f}{E k}.
\label{eq:alpha_F}
\end{equation}

Then
\begin{equation}
\partial_z(\nu u_z)-\alpha u=F.
\label{eq:reduced}
\end{equation}

Define the differential operator
\begin{equation}
L u
=
\partial_z(\nu u_z)-\alpha u.
\label{eq:L_def}
\end{equation}

The boundary conditions become
\begin{align}
\nu(0)u_z(0)&=\frac{\tau}{E k}, \label{eq:reduced_surface_bc}\\
u(-H)&=0. \label{eq:reduced_bottom_bc}
\end{align}

Thus the problem is
\begin{equation}
L u = F,
\end{equation}
with a bottom Dirichlet condition and a surface Neumann-flux condition.

\section{Change of variables}

Introduce
\begin{equation}
s=\exp\!\left(\frac{z}{h}\right).
\label{eq:s_def}
\end{equation}

Then
\[
z=h\log s,
\]
and the interval \([-H,0]\) maps to
\begin{equation}
s\in[s_b,s_s],
\qquad
s_b=e^{-H/h},
\qquad
s_s=1.
\end{equation}

The derivatives transform as
\begin{equation}
\frac{d}{dz}
=
\frac{s}{h}\frac{d}{ds}.
\end{equation}

Now rewrite the viscosity. Since
\[
\exp\!\left[-\frac{z+H}{h}\right]
=
e^{-H/h}s^{-1},
\qquad
\exp\!\left(\frac{z}{h}\right)=s,
\]
we have
\begin{equation}
\nu(z)
=
a+\frac{B}{s}+c s,
\qquad
B:=b e^{-H/h}.
\label{eq:nu_s}
\end{equation}

Therefore
\begin{equation}
s\nu(s)
=
c s^2+a s+B.
\end{equation}

Define the quadratic polynomial
\begin{equation}
P(s)
=
c s^2+a s+B.
\label{eq:P_def}
\end{equation}

Then
\begin{equation}
\nu(s)=\frac{P(s)}{s}.
\end{equation}

The homogeneous equation
\begin{equation}
\partial_z(\nu u_z)-\alpha u=0
\label{eq:hom_z}
\end{equation}
becomes
\begin{equation}
\frac{d}{ds}
\left(
P(s)\frac{dU}{ds}
\right)
-
\frac{\rho}{s}U=0,
\label{eq:hom_s_conservative}
\end{equation}
where
\begin{equation}
U(s)=u(z(s)),
\qquad
\rho=\alpha h^2
=
\frac{i f h^2}{E k}.
\label{eq:rho_def}
\end{equation}

Equivalently,
\begin{equation}
P(s)U''(s)+P'(s)U'(s)-\frac{\rho}{s}U(s)=0.
\label{eq:hom_s}
\end{equation}

This equation has regular singular points at
\[
s=0,
\]
and at the roots of \(P(s)\). It is a Fuchsian equation with four regular singular points, and therefore belongs to the general Heun class.

\section{Reduction to the general Heun equation}

Assume first that \(c\ne 0\) and that \(P(s)\) has two distinct roots \(r_1,r_2\):
\begin{equation}
P(s)=c(s-r_1)(s-r_2).
\label{eq:P_roots}
\end{equation}

Explicitly,
\begin{equation}
r_{1,2}
=
\frac{-a\pm\sqrt{a^2-4cB}}{2c},
\qquad
B=b e^{-H/h}.
\end{equation}

Introduce the dimensionless variable
\begin{equation}
x=\frac{s}{r_1},
\qquad
A=\frac{r_2}{r_1}.
\label{eq:x_A_def}
\end{equation}

Then the singular points are
\[
x=0,\qquad x=1,\qquad x=A,\qquad x=\infty.
\]

Writing
\[
U(s)=Y(x),
\]
equation \eqref{eq:hom_s} becomes
\begin{equation}
Y''+
\left[
\frac{1}{x-1}
+
\frac{1}{x-A}
\right]Y'
-
\frac{Q}{x(x-1)(x-A)}Y
=0,
\label{eq:heun_form}
\end{equation}
where
\begin{equation}
Q
=
\frac{\rho}{c r_1}
=
\frac{i f h^2}{E k\,c r_1}.
\label{eq:Q_def}
\end{equation}

The standard general Heun equation is
\begin{equation}
Y''
+
\left(
\frac{\gamma}{x}
+
\frac{\delta}{x-1}
+
\frac{\epsilon}{x-A}
\right)Y'
+
\frac{\alpha_H\beta_H x-q}
{x(x-1)(x-A)}Y
=0,
\label{eq:standard_heun}
\end{equation}
with the Fuchsian constraint
\begin{equation}
\alpha_H+\beta_H+1=\gamma+\delta+\epsilon.
\end{equation}

Comparing \eqref{eq:heun_form} and \eqref{eq:standard_heun}, one possible parameter choice is
\begin{equation}
\gamma=0,
\qquad
\delta=1,
\qquad
\epsilon=1,
\qquad
\alpha_H=0,
\qquad
\beta_H=1,
\qquad
q=Q.
\label{eq:heun_parameters}
\end{equation}

Therefore the homogeneous solutions may be represented in terms of general Heun functions.

A local solution about \(x=0\) may be written symbolically as
\begin{equation}
Y_1(x)
=
\operatorname{HeunG}
\left(
A,Q;0,1,0,1;x
\right),
\label{eq:heun_solution_1}
\end{equation}
with a second linearly independent solution obtained by the standard Frobenius construction. Because here \(\gamma=0\), the second solution may involve a logarithm. In practice it is often preferable to construct two independent Heun solutions about an ordinary point in the interval rather than about the singular point \(x=0\).

For the boundary-value problem below, we only require two independent homogeneous solutions. Denote them by
\begin{equation}
y_1(z),
\qquad
y_2(z),
\label{eq:y1_y2}
\end{equation}
where each is a solution of
\begin{equation}
L y=0.
\end{equation}

Equivalently,
\begin{equation}
y_j(z)
=
Y_j\!\left(
\frac{e^{z/h}}{r_1}
\right),
\qquad
j=1,2,
\end{equation}
where \(Y_j\) are independent solutions of \eqref{eq:heun_form}.

\section{Boundary-adapted homogeneous solutions}

Let \(\phi(z)\) be the homogeneous solution satisfying the bottom Dirichlet condition:
\begin{equation}
L\phi=0,
\qquad
\phi(-H)=0.
\label{eq:phi_def}
\end{equation}

Given two independent homogeneous solutions \(y_1,y_2\), define
\begin{equation}
\phi(z)
=
y_2(-H)y_1(z)-y_1(-H)y_2(z).
\label{eq:phi}
\end{equation}

Then
\[
\phi(-H)
=
y_2(-H)y_1(-H)-y_1(-H)y_2(-H)
=
0.
\]

Next let \(\psi(z)\) be the homogeneous solution satisfying the homogeneous surface-flux condition:
\begin{equation}
L\psi=0,
\qquad
\psi_z(0)=0.
\label{eq:psi_def}
\end{equation}

Since \(\nu(0)\ne 0\), this is equivalent to
\[
\nu(0)\psi_z(0)=0.
\]

Define
\begin{equation}
\psi(z)
=
y_2'(0)y_1(z)-y_1'(0)y_2(z).
\label{eq:psi}
\end{equation}

Then
\[
\psi_z(0)
=
y_2'(0)y_1'(0)-y_1'(0)y_2'(0)
=
0.
\]

Here primes denote differentiation with respect to \(z\).

\section{Weighted Wronskian}

For two homogeneous solutions \(\phi\) and \(\psi\), define the weighted Wronskian
\begin{equation}
W
=
\nu(z)
\left[
\phi(z)\psi_z(z)-\phi_z(z)\psi(z)
\right].
\label{eq:wronskian}
\end{equation}

This quantity is independent of \(z\). Indeed,
\begin{align}
\frac{dW}{dz}
&=
\frac{d}{dz}
\left\{
\nu
\left(
\phi\psi_z-\phi_z\psi
\right)
\right\} \notag\\
&=
\phi(\nu\psi_z)_z-\psi(\nu\phi_z)_z \notag\\
&=
\phi(\alpha\psi)-\psi(\alpha\phi) \notag\\
&=0.
\end{align}

\section{Green's function}

Define \(G(z,\zeta)\) by
\begin{equation}
L_z G(z,\zeta)=\delta(z-\zeta),
\label{eq:green_def}
\end{equation}
with homogeneous boundary conditions
\begin{align}
G(-H,\zeta)&=0, \label{eq:G_bottom}\\
G_z(0,\zeta)&=0. \label{eq:G_surface}
\end{align}

Then
\begin{equation}
G(z,\zeta)
=
\begin{cases}
\dfrac{\phi(z)\psi(\zeta)}{W},
& z<\zeta,\\[1.2em]
\dfrac{\phi(\zeta)\psi(z)}{W},
& z>\zeta.
\end{cases}
\label{eq:green}
\end{equation}

The Green's function is continuous at \(z=\zeta\):
\begin{equation}
G(\zeta^-,\zeta)=G(\zeta^+,\zeta),
\end{equation}
and its derivative has the jump
\begin{equation}
\nu(\zeta)
\left[
G_z(\zeta^+,\zeta)-G_z(\zeta^-,\zeta)
\right]
=
1.
\label{eq:jump}
\end{equation}

\section{Solution of the forced problem}

The inhomogeneous problem is
\begin{equation}
L u=F
=
\frac{\mathfrak f}{E k},
\label{eq:forced_problem}
\end{equation}
with
\begin{align}
u(-H)&=0,\\
\nu(0)u_z(0)&=\frac{\tau}{E k}.
\end{align}

The interior pressure-gradient forcing contributes
\begin{equation}
u_{\mathrm{int}}(z)
=
\int_{-H}^{0}
G(z,\zeta)F(\zeta)\,d\zeta
=
\frac{1}{E k}
\int_{-H}^{0}
G(z,\zeta)\mathfrak f(\zeta)\,d\zeta.
\label{eq:interior_solution}
\end{equation}

The nonhomogeneous surface stress contributes the boundary term
\begin{equation}
u_{\mathrm{surf}}(z)
=
-\frac{\tau}{E k}G(z,0).
\label{eq:surface_solution}
\end{equation}

Therefore the full solution is
\begin{equation}
\boxed{
u(z)
=
\frac{1}{E k}
\int_{-H}^{0}
G(z,\zeta)\mathfrak f(\zeta)\,d\zeta
-
\frac{\tau}{E k}G(z,0).
}
\label{eq:solution_z}
\end{equation}

Restoring the horizontal dependence,
\begin{equation}
\boxed{
u(x,y,z)
=
\frac{1}{E k}
\int_{-H}^{0}
G(z,\zeta)
\left[
p_x(x,y,\zeta)+i p_y(x,y,\zeta)
\right]\,d\zeta
-
\frac{\tau^x(x,y)+i\tau^y(x,y)}{E k}
G(z,0).
}
\label{eq:solution_full}
\end{equation}

\section{Summary}

For
\begin{equation}
\nu(z)
=
a+b\exp\!\left[-\frac{z+H}{h}\right]
+
c\exp\!\left(\frac{z}{h}\right),
\end{equation}
define
\begin{equation}
s=e^{z/h},
\qquad
B=b e^{-H/h},
\end{equation}
and
\begin{equation}
P(s)=c s^2+a s+B.
\end{equation}

The homogeneous equation becomes
\begin{equation}
\frac{d}{ds}
\left(
P(s)\frac{dU}{ds}
\right)
-
\frac{\rho}{s}U=0,
\qquad
\rho=\frac{i f h^2}{E k}.
\end{equation}

If \(c\neq 0\) and \(P(s)=c(s-r_1)(s-r_2)\), define
\begin{equation}
x=\frac{s}{r_1},
\qquad
A=\frac{r_2}{r_1},
\qquad
Q=\frac{\rho}{c r_1}.
\end{equation}

The homogeneous equation becomes
\begin{equation}
Y''
+
\left[
\frac{1}{x-1}
+
\frac{1}{x-A}
\right]Y'
-
\frac{Q}{x(x-1)(x-A)}Y
=0,
\end{equation}
which is a general Heun equation with parameters
\begin{equation}
\gamma=0,
\qquad
\delta=1,
\qquad
\epsilon=1,
\qquad
\alpha_H=0,
\qquad
\beta_H=1,
\qquad
q=Q.
\end{equation}

Let \(y_1,y_2\) be any two independent homogeneous solutions. Define
\begin{align}
\phi(z)
&=
y_2(-H)y_1(z)-y_1(-H)y_2(z),
\\[0.5em]
\psi(z)
&=
y_2'(0)y_1(z)-y_1'(0)y_2(z).
\end{align}

Then
\begin{equation}
W
=
\nu(z)
\left[
\phi(z)\psi_z(z)-\phi_z(z)\psi(z)
\right]
\end{equation}
is constant, and
\begin{equation}
G(z,\zeta)
=
\begin{cases}
\dfrac{\phi(z)\psi(\zeta)}{W},
& z<\zeta,\\[1.2em]
\dfrac{\phi(\zeta)\psi(z)}{W},
& z>\zeta.
\end{cases}
\end{equation}

The velocity is
\begin{equation}
\boxed{
u(x,y,z)
=
\frac{1}{E k}
\int_{-H}^{0}
G(z,\zeta)
\left[
p_x(x,y,\zeta)+i p_y(x,y,\zeta)
\right]\,d\zeta
-
\frac{\tau^x(x,y)+i\tau^y(x,y)}{E k}G(z,0).
}
\end{equation}

\section{Special cases}

\subsection{Case \(c=0\)}

If \(c=0\), then
\[
\nu(z)=a+b\exp\!\left[-\frac{z+H}{h}\right].
\]
This is the single-exponential case. The homogeneous equation reduces to a Gauss hypergeometric equation. With
\[
X(z)
=
-\frac{b}{a}
\exp\!\left[-\frac{z+H}{h}\right],
\]
and
\[
\lambda
=
h\sqrt{\frac{i f}{E k\,a}},
\]
a homogeneous basis is
\begin{align}
y_+(z)
&=
X(z)^\lambda
{}_2F_1
\left(
\lambda,\lambda+1;2\lambda+1;X(z)
\right),
\\[0.5em]
y_-(z)
&=
X(z)^{-\lambda}
{}_2F_1
\left(
-\lambda,1-\lambda;1-2\lambda;X(z)
\right).
\end{align}

\subsection{Case \(b=0\)}

If \(b=0\), then
\[
\nu(z)=a+c\exp\!\left(\frac{z}{h}\right).
\]
This is also reducible to a Gauss hypergeometric equation by using
\[
X(z)
=
-\frac{c}{a}e^{z/h}.
\]

\subsection{Constant-viscosity limit}

If \(b=c=0\), then
\[
\nu(z)=a,
\]
and the equation becomes
\begin{equation}
a u_{zz}-\alpha u=F.
\end{equation}

Define
\begin{equation}
m=\sqrt{\frac{\alpha}{a}}
=
\sqrt{\frac{i f}{E k\,a}}.
\end{equation}

The homogeneous solution is
\begin{equation}
u_h(z)
=
C_1\cosh\!\left[m(z+H)\right]
+
C_2\sinh\!\left[m(z+H)\right].
\end{equation}

A bottom-Dirichlet solution is
\begin{equation}
\phi(z)=\sinh\!\left[m(z+H)\right],
\end{equation}
and a homogeneous surface-flux solution is
\begin{equation}
\psi(z)=\cosh(mz),
\end{equation}
since
\[
\psi_z(0)=0.
\]

This gives the usual finite-depth Ekman solution with bottom no-slip and imposed surface stress.

\section{Remarks}

\begin{itemize}
    \item The two-exponential viscosity profile
    \[
    \nu(z)=a+b e^{-(z+H)/h}+c e^{z/h}
    \]
    leads to a general Heun equation, whereas the one-exponential profile leads to a Gauss hypergeometric equation.

    \item The Green's-function representation does not depend on the particular normalization of \(y_1,y_2\). Any two independent homogeneous solutions may be used.

    \item Since
    \[
    \alpha=\frac{i f}{E k},
    \]
    the homogeneous solutions are generally complex-valued.

    \item The branch choices in the Heun functions must be made consistently if the singularities lie near the physical interval.

    \item The physical interval in the \(s\)-variable is
    \[
    e^{-H/h}\le s\le 1.
    \]
    The solution is regular provided \(P(s)\neq 0\) on this interval, which is equivalent to \(\nu(z)>0\).

    \item If a root of \(P(s)\) lies on the physical interval, then the viscosity vanishes and the boundary-value problem becomes singular.
\end{itemize}

\begin{thebibliography}{99}

\bibitem{Ekman1905}
V. W. Ekman.
\newblock On the influence of the Earth's rotation on ocean-currents.
\newblock \emph{Arkiv f\"or Matematik, Astronomi och Fysik}, 2(11):1--52, 1905.

\bibitem{Gill1982}
A. E. Gill.
\newblock \emph{Atmosphere--Ocean Dynamics}.
\newblock Academic Press, 1982.

\bibitem{Pedlosky1987}
J. Pedlosky.
\newblock \emph{Geophysical Fluid Dynamics}.
\newblock 2nd edition, Springer, 1987.

\bibitem{Vallis2017}
G. K. Vallis.
\newblock \emph{Atmospheric and Oceanic Fluid Dynamics}.
\newblock 2nd edition, Cambridge University Press, 2017.

\bibitem{Heun1889}
K. Heun.
\newblock Zur Theorie der Riemann'schen Functionen zweiter Ordnung mit vier Verzweigungspunkten.
\newblock \emph{Mathematische Annalen}, 33:161--179, 1889.

\bibitem{Ronveaux1995}
A. Ronveaux, editor.
\newblock \emph{Heun's Differential Equations}.
\newblock Oxford University Press, 1995.

\bibitem{SlavyanovLay1999}
S. Yu. Slavyanov and W. Lay.
\newblock \emph{Special Functions: A Unified Theory Based on Singularities}.
\newblock Oxford University Press, 1999.

\bibitem{Maier2007}
R. S. Maier.
\newblock The 192 solutions of the Heun equation.
\newblock \emph{Mathematics of Computation}, 76(258):811--843, 2007.

\bibitem{DLMFHeun}
NIST Digital Library of Mathematical Functions.
\newblock \emph{Heun Functions}.
\newblock Chapter 31.
\newblock \url{https://dlmf.nist.gov/31}

\bibitem{DLMFHyper}
NIST Digital Library of Mathematical Functions.
\newblock \emph{Gauss Hypergeometric Function}.
\newblock Chapter 15.
\newblock \url{https://dlmf.nist.gov/15}

\bibitem{Olver2010}
F. W. J. Olver, D. W. Lozier, R. F. Boisvert, and C. W. Clark, editors.
\newblock \emph{NIST Handbook of Mathematical Functions}.
\newblock Cambridge University Press, 2010.

\bibitem{Zettl}
A. Zettl.
\newblock \emph{Sturm--Liouville Theory}.
\newblock American Mathematical Society, 2005.

\bibitem{CoddingtonLevinson}
E. A. Coddington and N. Levinson.
\newblock \emph{Theory of Ordinary Differential Equations}.
\newblock McGraw--Hill, 1955.

\bibitem{ArfkenWeber}
G. B. Arfken, H. J. Weber, and F. E. Harris.
\newblock \emph{Mathematical Methods for Physicists}.
\newblock 7th edition, Academic Press, 2013.

\end{thebibliography}

\end{document}